\section{Tool-Vergleich: Celonis vs. Python}
\label{sec:toolvergleich}

% Umfang: 1-2 Seiten
% Assignment: "Pick two out of three ... Compare the results!"
% - Gegenüberstellung der beiden gewählten Werkzeuge (Vergleichstabelle):
%   Bedienung, Exploration/Visualisierung, Discovery, Conformance (Happy Path vs.
%   Alignments auf normativem Modell), Reproduzierbarkeit, ML-Erweiterbarkeit, Kosten.
% - Stärken Celonis (schnelle visuelle Exploration, Dashboarding, PQL) vs.
%   Python/PM4Py (Flexibilität, spezifische Algorithmen) \parencite{berti_process_2019,vogelgesang_celonis_2022}
% - Ergebnisvergleich: Abweichungen bei Varianten, Durchlaufzeiten, Konformitätsquote?
%   * Durchlaufzeit International: Celonis zeigt 86 d (alle Cases, Throughput Times),
%     75 d (ohne Extremwerte, Process Overview) und 66 d (Happy Path). Ein mean() in
%     Python ergibt ~86 d -> Beispiel dafür, dass Standard-KPIs in BI-Tools still
%     trimmen und von Rohdaten-Berechnungen abweichen (vgl. Abschnitt 3.1).
%   * Konformität: Celonis prüft gegen ein aus den Top-Varianten generiertes Modell
%     (81 % / 38 %), PM4Py gegen normatives Modell (Alignments) -> Werte nicht 1:1 vergleichbar.
%   * Datenqualität: 134 Domestic-Cases mit nur "Declaration SAVED" (0 d Durchlaufzeit)
%     -> beeinflussen Mittelwerte je nach Filterung in beiden Tools unterschiedlich.

\todo{Tool-Vergleich (gemeinsam)}
